\documentclass[11pt]{article}
\usepackage[margin=1in]{geometry}
\usepackage{amsmath,amssymb,amsthm,mathtools,microtype,hyperref}
\hypersetup{hidelinks}
\newtheorem{theorem}{Theorem}
\newtheorem{lemma}{Lemma}
\newtheorem{proposition}{Proposition}
\newcommand{\ddc}{dd^c}
\newcommand{\PSH}{\operatorname{PSH}}
\newcommand{\MA}{\operatorname{MA}}
\newcommand{\ext}{\mathrm{ext}}
\newcommand{\PP}{\mathbb P}
\title{Nonanalytic imposed poles and moving volume data on a ruled surface}
\author{}
\date{October 8, 2026}
\begin{document}
\maketitle
\noindent\textbf{Checked restricted construction.} The stated results have
passed independent review, statement-only reconstruction and final comparison.
The unrestricted S18 problem remains open.

\begin{abstract}
We give a concrete moving-volume construction on $\PP^1\times\PP^1$
with a positive genuinely nonclosed background and an imposed unbounded
nonanalytic ordinary singularity type. The upper and lower masses optimize
over every potential of that type, including rough fiber-dependent potentials
and references with infinite first moment. The calculation uses closed
ordinary-order comparison in an auxiliary background and requires only
local integrability in its nonclosed correction.
They are distinct from the fixed-density Gauduchon variational volume, whose
full-test infimum and subtraction cost admit explicit formulas. A smooth nef
endpoint gives logarithmic variational degeneration while the upper singularity
mass stays positive. Exact smooth-gauge and pure-pullback laws distinguish
these data from exceptional-curvature subtraction and from arbitrary changes
of the reference density. This is a restricted construction; the general
nonclosed moving-product programme is not resolved.
\end{abstract}

\section{Definitions and results}
Use $\ddc=(i/(2\pi))\partial\bar\partial$ and
$d^c=(i/(4\pi))(\bar\partial-\partial)$. Let $w$ be the Fubini--Study
probability form on $\PP^1$, and write $w_1,w_2$ for its pullbacks to
$X=\PP^1\times\PP^1$. In an affine coordinate,
\[
 w=\frac{i\,dz\wedge d\bar z}{2\pi(1+|z|^2)^2},\qquad
 G_0=\log\frac{|z|^2}{1+|z|^2},\qquad \ddc G_0=[0]-w.
\]
Unitary transport gives $G_p$ for any $p\in\PP^1$, normalized by
$\max G_p=0$. Fix $0<\alpha<1$ and set
\[
 U=1-(1-G_p)^\alpha,\qquad u=\operatorname{pr}_1^*U,
 \qquad D=\{p\}\times\PP^1.
\]
All potentials use their genuine canonical upper representatives. We write
$v\sim u$ for bounded difference, with the inequalities retained at common
poles. For $c\geq\alpha$ and a smooth function $a>0$ on the first factor,
put $\theta=cw_1+aw_2$ and define
\[
 \sigma_-(\theta,u)=\inf_{v\in\PSH(X,\theta),\ v\sim u}
              \MA^{\ext}_\theta(v)(X),\qquad
 \sigma_+(\theta,u)=\sup_{v\in\PSH(X,\theta),\ v\sim u}
              \MA^{\ext}_\theta(v)(X).
\]
The products are the finite-level non-pluripolar products, expanded against
the smooth background. These extrema concern the entire ordinary type. Together with the separately
defined fixed-density variational volume and its minimizing-test cost, they
are the moving data computed here. The imposed ordinary type is retained
as part of these data: its trivial ideals do not identify it with the
bounded type.

\begin{theorem}\label{thm:mass}
The reference $u$ is $\theta$-psh and has an unbounded nonanalytic pole on
$D$, zero Lelong numbers, and $\mathcal I(tu)=\mathcal O_X$ for every $t>0$.
Every $v\in\PSH(X,\theta)$ with $v\sim u$ has positive finite mass and
\begin{align}
 \MA^{\ext}_\theta(v)(X)
     &=\int_X\theta^2+2\int_Xv\,\ddc\theta,\label{eq:linear-mass}\\
 \sigma_-(\theta,u)&=2c\min a,\qquad
 \sigma_+(\theta,u)=2c\max a.\label{eq:extrema}
\end{align}
Both bounds are approached by actual representatives of that exact type.
When $a$ is nonconstant the smooth positive background is nonclosed.
\end{theorem}

For a smooth positive $m$ on the first factor with $\int m w=1$, let
$\mu=mw_1w_2$. For a smooth positive Gauduchon test $T$ of bidegree $(1,1)$,
so that $\ddc T=0$, define
\[
 D_\mu(T)=\int_X\sqrt{T^2/\mu}\,\mu,\qquad
 \lambda_\mu(\theta)=\inf_T\frac{\int_X\theta\wedge T}{D_\mu(T)},
 \qquad V_\mu(\theta)=\lambda_\mu(\theta)^2.
\]
Every smooth positive Gauduchon test is included in this infimum.

\begin{theorem}\label{thm:dual}
Writing $H=\int_{\PP^1}(m/a)w$, one has
\[
 V_\mu(\theta)=\frac{2c}{H}.
\]
The infimum is attained by the smooth positive closed test
\[
 T=\frac ga w_1+w_2,\qquad g=\frac{V_\mu(\theta)m}{2a}.
\]
The mass extrema in \eqref{eq:extrema} are independent of $m$, whereas
this fixed-density volume can depend on $m$ when $a$ is nonconstant.
\end{theorem}

The full-type mass calculation, the all-test variational formula and the
pure-pullback law agree with the checked ruled construction
\cite{Ruled}, specialized to zero logarithmic divisor coefficient.
They are proved here to fix their actual application to the moving
segment and its limiting cost. Existing ordinary-type volumes and their
pure-modification laws receive their credit in \cite{BGL}.
No literature-priority claim is made.

\section{The imposed type and small-background representatives}
Let $\chi(t)=1-(1-t)^\alpha$ for $t\leq0$. Then
\[
 \chi'(t)=\alpha(1-t)^{\alpha-1},\qquad
 \chi''(t)=\alpha(1-\alpha)(1-t)^{\alpha-2}.
\]
On the punctured curve,
$\ddc U=-\chi'(G_p)w+\chi''(G_p)dG_p\wedge d^cG_p\geq-\alpha w$.
After adding a local potential of $\alpha w$, this is a locally bounded-above
psh function on the punctured disk. Its genuine psh extension has the same
lower curvature bound. Its logarithmic coefficient is zero, since
$U/\log|z-p|\to0$. In particular the positive curve current
$Cw+\ddc U$ has no atom at $p$ for $C\geq1$.
The same local calculation proves $u\in\PSH(X,\theta)$.
Moreover $e^{-tU}$ grows more slowly than every fixed negative power of
$|z-p|$, and hence is locally integrable for every $t>0$. All holomorphic
germs are therefore in $\mathcal I(tu)$. An unbounded analytic singularity
on this smooth divisor has a positive logarithmic coefficient; the imposed
type is thus genuinely nonanalytic.

\begin{lemma}\label{lem:small}
There are $U_L\sim U$ with $\sup U_L=0$,
$\ddc U_L\geq-\varepsilon_Lw$, $\varepsilon_L\to0$, and
$U_L\to0$ in $L^1(\PP^1)$.
\end{lemma}
\begin{proof}
For $L>0$ put $\varepsilon_L=\chi'(-L)$ and define
\[
 \chi_L(t)=
 \begin{cases}
   \chi(t)+C_L,&t<-L,\\
   \varepsilon_Lt,&-L\leq t\leq0,
 \end{cases}\qquad
 C_L=-\varepsilon_LL-\chi(-L).
\]
The two values and first derivatives match. This is convex with
$0\leq\chi_L'\leq\varepsilon_L$. Set $U_L=\chi_L(G_p)$.
The preceding local extension argument, or smooth convex approximation
followed by local decreasing convergence, gives its genuine curvature bound.
Near $p$ the difference $U_L-U$ is the finite constant $C_L$; on the remaining
compact part it is bounded. Thus $U_L\sim U$, including their common pole.
Here $C_L\geq0$: with $x=1+L$, the function
$(1-\alpha)x^\alpha+\alpha x^{\alpha-1}-1$ vanishes at $x=1$ and has
nonnegative derivative for $x\geq1$. On $[-L,0]$, convexity gives
$\chi(t)\leq\chi'(-L)t$. Consequently $U\leq U_L\leq0$.
For each point other than $p$, eventually $U_L=\varepsilon_LG_p\to0$.
Since $U\in L^1$, dominated convergence proves the stated $L^1$ limit.
\end{proof}

This supplies actual small-background representatives without erasing the
pole by a maximum with a bounded function. If a smooth $g$ satisfies
$cw+\ddc g\geq\delta w$, then $g+U_L$ is $cw$-psh for
$\varepsilon_L<\delta$ and remains of the ordinary type $U$.

\section{All rough competitors and the mass extrema}
We use the closed ordinary-singularity comparison of
\cite[Theorem~1.16]{BEGZ}. Its hypotheses are compact K\"ahler ambient,
closed positive currents in common cohomology classes, and small unbounded
locus, meaning local boundedness outside a closed locally complete
pluripolar set. We check each hypothesis and use the comparison in both
ordinary orders and in degrees one and two. No nonclosed comparison
or finite-energy theorem is invoked.

Choose $K=C(w_1+w_2)$, with $C$ sufficiently large that $K\geq\theta$
and $C\geq1$. Write $S_0=K+\ddc u$. Its actual closed products are
\[
 \langle S_0\rangle=S_0,\qquad
 \langle S_0^2\rangle=2C(Cw_1+\ddc u)w_2,
 \qquad\int_X\langle S_0^2\rangle=\int_XK^2=2C^2.
\]
Indeed the curve current has total mass $C$ and no atom at $p$.
On $X\setminus D$ its pulled-back square is zero by dimension.
The displayed positive finite product has zero mass on $D$ and is
exactly the non-pluripolar extension of the bounded local product.

Take any $v\in\PSH(X,\theta)$ with $v\sim u$, and put
$S=K+\ddc v$. Both $S,S_0$ are closed positive currents in the
K\"ahler class $\{K\}$ on the compact K\"ahler surface $X$.
Their potentials are locally bounded off the complete pluripolar divisor
$D$, and have mutual bounded order. The cited theorem, applied in both
orders to a single current and then to the repeated two-current tuples,
gives
\[
 \{\langle S\rangle\}=\{\langle S_0\rangle\}=\{K\},\qquad
 \int_X\langle S^2\rangle=\int_X\langle S_0^2\rangle=\int_XK^2.
\]
Here $\langle S\rangle$ is $S$ restricted off $D$ and extended
by zero. Thus $S-\langle S\rangle=1_DS$ is positive and its trace is
\[
 \int_X(S-\langle S\rangle)\wedge K=0.
\]
A positive current of zero strictly positive trace is zero, as follows
from the local Hermitian matrix inequality for its measure coefficients.
Consequently $\langle S\rangle=S$, including absence of polar trace.
This proves the required first-current identity for every rough competitor;
it does not infer it from vanishing generic divisorial Lelong numbers alone.

Set $R=\theta-K$. The finite-level smooth-background expansion now gives
\[
 \MA_\theta^{\ext}(v)=\langle S^2\rangle+2R\wedge S+R^2.
\]
All terms are finite signed measures, and the left side is positive.
Integrating, using full closed mass and distributional Stokes, gives
\[
 \int K^2+2\int R\wedge K+\int R^2+2\int R\wedge\ddc v
     =\int\theta^2+2\int v\,\ddc\theta.
\]
This proves \eqref{eq:linear-mass}, including for every rough competitor.

Average $v$ over the compact unitary action on the second factor.
The background is invariant. Averaging its current inequality and using
canonical psh representatives gives a genuine invariant potential
$\operatorname{pr}_1^*q$, with $q\sim U$. For completeness, local addition
of a smooth sufficiently positive quadratic makes all translates psh;
their Haar average is locally integrable by Fubini, satisfies the submean
inequality, and has its canonical upper representative. The bounded
inequalities with $u$ pass through the average and then through canonical
regularization. Off $D$ the average is finite, and on $D$ it equals
$-\infty$. Positivity says $cw+\ddc q\geq0$.

Since $\ddc\theta=\ddc a\wedge w_2$, the functional in
\eqref{eq:linear-mass} is unchanged by this averaging, and
\[
 \MA_\theta^{\ext}(v)(X)=2\int_{\PP^1}a(cw+\ddc q).
\]
The measure in parentheses is positive of total mass $c$. This proves both
global bounds and positive finiteness.

To approach either bound, choose a point $x$ attaining the corresponding
extremum of $a$. There are smooth strictly positive curve measures
$\zeta_j$ of total mass $c$ converging weakly to $c\delta_x$; for instance,
combine a smooth approximate identity near $x$ with a small positive
multiple of $w$. Solve $\ddc g_j=\zeta_j-cw$ on the sphere. The smooth
solution exists because the right side has zero integral, either by the
explicit Green kernel or by the zero-mean scalar Poisson equation.
For each $j$ choose $L_j$ so large that
$\varepsilon_{L_j}$ is smaller than half the positive lower bound of
$\zeta_j/w$ and $|\int U_{L_j}\ddc a|<1/j$.
Then $g_j+U_{L_j}$ is an actual $cw$-psh representative of $U$'s type and
\[
 \int a(cw+\ddc(g_j+U_{L_j}))\longrightarrow c a(x).
\]
Pulling back proves both exact extrema in Theorem~\ref{thm:mass}.


\section{An exact first-moment boundary}
The preceding proof holds for every $0<\alpha<1$. For the same closed
reference $K=C(w_1+w_2)$, its first moment has a separate threshold:
\[
 \int_X(-u)\langle (K+\ddc u)^2\rangle<\infty
       \quad\Longleftrightarrow\quad\alpha<\tfrac12.
\]
On a punctured disk put $t=-G_p$. As $t\to\infty$, $-U$ is comparable
to $t^\alpha$ and $\chi''(G_p)$ to $t^{\alpha-2}$.
The angular integral of $dG_p\wedge d^cG_p$ is comparable to $dt$.
The singular contribution to the actual curve moment is therefore
comparable to $\int^\infty t^{2\alpha-2}\,dt$.
The remaining smooth-volume contribution is finite by exponential radial
area decay. The moment is infinite for $1/2\leq\alpha<1$, logarithmically
at equality. This is an actual separated non-pluripolar measure calculation,
not an extra finiteness hypothesis in the nonclosed mass formula.
In particular $\int v\,\ddc\theta$ uses only the local integrability
of $v$ against a smooth form, throughout this enlarged parameter range.

\section{The full Gauduchon formula and its moving cost}
Let $H=\int(m/a)w$, $V=2c/H$ and $g=Vm/(2a)$, so $\int gw=c$.
Choose a smooth $b$ with $\ddc b=(g-c)w$. Then
\[
 \Omega=\theta+\ddc\operatorname{pr}_1^*b=gw_1+aw_2>0,
 \qquad\Omega^2=V\mu.
\]
For every smooth positive Gauduchon test $S$, Stokes and the pointwise
two-dimensional arithmetic--geometric mean inequality give
\[
 \int\theta\wedge S=\int\Omega\wedge S
       \geq\sqrt V\int\sqrt{S^2/\mu}\,\mu=\sqrt V D_\mu(S).
\]
The test $T=\Omega/a=(g/a)w_1+w_2$ is positive and closed, because its
first coefficient depends only on the first curve. Moreover
\[
 \int\theta T=2c,\qquad T^2=\frac{Vm}{a^2}w_1w_2,
 \qquad D_\mu(T)=\sqrt V H.
\]
It attains the lower bound, proving Theorem~\ref{thm:dual} for all tests.

For $c_s=c-sd\geq\alpha$, $a_s=a-sb>0$ on a parameter interval, put
$\theta_s=c_sw_1+a_sw_2$ and $\eta=dw_1+bw_2$. Keep $m$ fixed, and
use the preceding notation with subscript $s$. Direct differentiation gives
\[
 H_s'=\int\frac{bm}{a_s^2}w,\qquad
 \int\eta T_s=d+\frac{V_s}{2}H_s',\qquad
 \frac{d\sqrt{V_s}}{ds}
       =-\frac{\int\eta T_s}{D_\mu(T_s)}.
\]
This is the actual full-test subtraction derivative. Its integral on any
compact positive segment equals the difference of the two endpoint root
volumes. It is separate from the upper and lower singularity masses.

\section{A logarithmic nef endpoint}
Let $c=1$, $m=1$, and
$\tau=(|z|^2-1)/(|z|^2+1)$ on the first sphere. Its distribution under
$w$ is $d\tau/2$ on $[-1,1]$: radial integration gives
$dt/(1+t)^2=d\tau/2$, with $t=|z|^2$. Take $a=2+\tau$ and subtract
$\eta=w_2$. For $0\leq s<1$,
\[
 H_s=\frac12\log\frac{3-s}{1-s},\qquad
 V_s=\frac4{\log((3-s)/(1-s))},\qquad
 \sigma_-(s)=2(1-s),\quad\sigma_+(s)=2(3-s).
\]
At $s=0$ the raw zero-potential mass is $2\int a w=4$, while the lower
mass, upper mass and fixed-density volume are $2$, $6$ and $4/\log3$.
All these backgrounds are genuinely nonclosed. As $s\uparrow1$ the root
volume tends to zero and its derivative is
\[
 -\frac{2}{(1-s)(3-s)[\log((3-s)/(1-s))]^{3/2}}.
\]
The improper integral of the cost is finite, equal to $2/\sqrt{\log3}$.
The upper singularity mass instead tends to $4$.

The endpoint $\theta_1=w_1+(1+\tau)w_2$ is smooth semipositive and nef,
since adding any positive multiple of $w_1+w_2$ makes it positive. Its
all-test root-volume infimum is zero, by nonnegativity and comparison with
$\theta_s$. It is not big. Indeed a strict current would dominate
$\varepsilon(w_1+w_2)$ after changing $\varepsilon>0$. Averaging its
potential over the second sphere preserves that strict inequality and
produces a genuine base potential. Its vertical coefficient is still
$1+\tau$, which would therefore have to be at least $\varepsilon$
almost everywhere. This is impossible near its zero fiber.
The endpoint thus lies outside the big segment; it is not a counterexample
to any asserted equivalence between bigness and positive variational volume.

\section{Smooth gauges and pure pullback}
For smooth real $f$, set $\theta'=\theta+\ddc f$, $u'=u-f$ and $v'=v-f$.
The actual currents coincide, so the finite-level local construction
transports the extended measure and both mass extrema. Against every
Gauduchon test the numerator also coincides by Stokes. A parameter-dependent
smooth gauge contributes $\ddc\dot f_s$ to the derivative and its pairing
with $T_s$ is zero.

\begin{proposition}\label{prop:mod}
For a proper modification $\pi:Y\to X$ between smooth connected compact
complex manifolds and a smooth real background $\theta$, pullback is a
canonical bijection of the genuine $\theta$-psh functions with the genuine
$\pi^*\theta$-psh functions. Ordinary types correspond. For these
potentials,
\[
 \pi_*\MA_{\pi^*\theta}^{\ext}(\pi^*v)=\MA_\theta^{\ext}(v).
\]
Thus the mass extrema are preserved for the pure pullback background and
type whenever the class of competitors is nonempty.
\end{proposition}
\begin{proof}
Let $Z$ be an analytic center containing the non-isomorphism locus.
On a relatively compact coordinate ball choose smooth $\rho$ with
$\ddc\rho\geq\theta$. If $b$ is $\pi^*\theta$-psh, then
$b+\pi^*\rho$ is psh upstairs. Off $Z$ it descends to a psh function;
properness and local upper bounds upstairs give a local upper bound near
$Z$. The removable-singularity theorem \cite[I.5.23--5.24]{Demailly}
gives the unique canonical psh extension. Subtracting $\rho$ gives a genuine
quasi-psh function downstairs. Its sharper $\theta$-positivity holds off
$Z$ and also on $Z$: the added current on $Z$ is positive, whereas the
smooth correction has no mass there. Pullback agrees upstairs off the
exceptional analytic set, hence almost everywhere, and the two genuine
quasi-psh representatives agree canonically. The converse follows from
local psh pullback and the current bound.

Bounded inequalities descend off $Z$ and hold everywhere by canonical
recovery of quasi-psh functions from their almost-everywhere values. Thus
ordinary equivalence is preserved in both directions. Off the exceptional
set all finite-level local products transform by the biholomorphism.
Each bounded local product, and then each non-pluripolar product, gives
zero mass to analytic pluripolar sets; this remains true for the
smooth-background expansion. Both extended measures therefore give zero
mass to the respective excluded analytic sets. The equality on their
complements proves the measure identity, even when total mass is infinite.
Taking the full infimum and supremum over the bijective type proves the
last assertion.
\end{proof}

For a point blowup the pullback backgrounds in the positive segment are
also big. One direct justification uses the exceptional divisor $E$.
Choose a smooth curvature $\beta_E$ of $\mathcal O(E)$ whose restriction
to $E$ is the negative Fubini--Study form. For a closed K\"ahler form
$K_0=w_1+w_2$, the form $\pi^*K_0-t\beta_E$ is K\"ahler for sufficiently
small $t>0$. Here is the local positivity check. In finitely many blowup
charts $(x,y)\mapsto(x,xy)$ with $y$ in compact sets, the matrix $P$ of
$\pi^*K_0$ satisfies $P_{x\bar x}\geq c_0>0$,
$|P_{x\bar y}|\leq C|x|$, and its Schur complement is at least
$c_1|x|^2$. For $B=-\beta_E$, shrink a fixed chart neighborhood so that
$B_{y\bar y}\geq b_0>0$, with all entries bounded. The transverse
coefficient of $P+tB$ is at least $c_0/2$ for small $t$, and its Schur
complement is at least
\[
 c_1|x|^2+t b_0-C'\bigl(t|x|+t|x|^2+t^2\bigr).
\]
First fix the neighborhood small enough, then $t$ small enough, to make
the last three terms at most $t b_0/2$. This proves strict positivity on
that fixed neighborhood. On its compact complement, strict positivity
of $\pi^*K_0$ persists under the bounded perturbation. The metric with
the required restriction to $E$ exists because
$\mathcal O(E)|_E=\mathcal O_{\PP^1}(-1)$: extend its negative
Fubini--Study metric smoothly to the line bundle. If $\theta_s\geq a_0K_0$ and
$\ddc\ell_E=[E]-\beta_E$, then
\[
 \pi^*\theta_s+\ddc(a_0t\ell_E)
   \geq a_0(\pi^*K_0-t\beta_E)+a_0t[E]
\]
is an actual strict current. This gives a nonclosed big example upstairs.
Subtracting exceptional curvature is a different background, and is not
covered by Proposition~\ref{prop:mod}'s equality. Nor do we identify an
arbitrary smooth positive upstairs probability with $\pi^*\mu$, which
degenerates on $E$.

The construction does not assert a general nonclosed positive-product
theory, unrestricted birational equality for changed backgrounds,
density-independent Gauduchon volume, or a solution of the general
nef-endpoint bridge. The closed comparison in \cite{BEGZ} supplies only
the precisely checked auxiliary closed foundations. The nonclosed
optimization, full-test formula and limiting subtraction cost have
their stated direct proofs. General S18 remains open.

\begin{thebibliography}{9}
\bibitem{BEGZ} S.~Boucksom, P.~Eyssidieux, V.~Guedj and A.~Zeriahi,
\emph{Monge--Amp\`ere equations in big cohomology classes},
Acta Math. \textbf{205} (2010), 199--262.
\href{https://www.math.univ-toulouse.fr/~guedj/fichierspdf/BEGZ10.pdf}{Author-hosted published primary text}.
\bibitem{Demailly} J.-P.~Demailly, \emph{Complex Analytic and Differential
Geometry}, Chapter~I, Theorems~5.23--5.24.
\href{https://www-fourier.univ-grenoble-alpes.fr/~demailly/manuscripts/agbook.pdf}{University-hosted primary text}.
\bibitem{BGL} S.~Boucksom, V.~Guedj and C.~H.~Lu,
\emph{Volumes of Bott--Chern classes}, Peking Math.~J. (2025),
Definition~3.1 and Proposition~3.3.
\href{https://d-nb.info/1374372846/34}{Published primary text}.
\bibitem{Ruled} \emph{Restricted masses and ordinary strict-current types
on a nonclosed ruled surface}, checked restricted construction, October~2026,
Sections~2 and~5.
\href{https://github.com/MingchenXia/non-kahler-pluripotential-theory/blob/2ec45dec6172ced6e2af18b2b5fe657c07aeee9e/research/papers/S11-nonclosed-ruled-restricted-masses-and-strict-types.tex}{Retained checked source}.
\end{thebibliography}
\end{document}
